\chapter*{ПРИЛОЖЕНИЕ Б. ТОПОЛОГИЯ РАЗВЕРТЫВАНИЯ И ПРОФИЛИ ЗАПУСКА}
\phantomsection
\addcontentsline{toc}{chapter}{ПРИЛОЖЕНИЕ Б. ТОПОЛОГИЯ РАЗВЕРТЫВАНИЯ И ПРОФИЛИ ЗАПУСКА}

\section*{Б.1 Минимальный compose-контур}
\phantomsection
\addcontentsline{toc}{section}{Б.1 Минимальный compose-контур}

\begin{longtable}{@{}p{0.3\textwidth}p{0.64\textwidth}@{}}
\texttt{ui} & SPA-интерфейс оператора \\
\texttt{api} & Управление конфигурацией, событиями, аутентификацией, RBAC и служебными маршрутами \texttt{/health}, \texttt{/metrics} \\
\texttt{analytics-worker} & Захват потоков, инференс, сопровождение объектов, правила, таблица исходящих сообщений \\
\texttt{integration-worker} & Публикация и доставка событий, повторы, очередь недоставленных сообщений \\
\texttt{postgres} & Хранение конфигурации, событий, таблицы исходящих сообщений и журнала доставок \\
\texttt{rabbitmq} & Асинхронный транспорт событий \\
\texttt{prometheus} & Сбор метрик \\
\texttt{grafana} & Визуализация метрик \\
\texttt{webhook-mock} & Тестовый HTTP-получатель \\
\end{longtable}

\section*{Б.2 Профили запуска}
\phantomsection
\addcontentsline{toc}{section}{Б.2 Профили запуска}

CPU-профиль разработки:
\begin{itemize}
    \item 1--3 видеопотока;
    \item базовая легкая YOLO-модель на CPU;
    \item акцент на функциональные и интеграционные проверки.
\end{itemize}

Профиль дальнейшего аппаратного ускорения:
\begin{itemize}
    \item 1--10 видеопотоков;
    \item инференс с аппаратным ускорением;
    \item фиксация p95-задержек и устойчивости доставки событий в отдельной серии прогонов.
\end{itemize}

\section*{Б.3 Операционные команды}
\phantomsection
\addcontentsline{toc}{section}{Б.3 Операционные команды}

\begin{verbatim}
# запуск стека
make up-d

# просмотр статуса
make ps

# применение миграций
make migrate-up

# запуск smoke-сценария
make smoke-e2e

# остановка стека
make down
\end{verbatim}

\section*{Б.4 Рекомендации по воспроизводимости}
\phantomsection
\addcontentsline{toc}{section}{Б.4 Рекомендации по воспроизводимости}

\begin{enumerate}
    \item Фиксировать версии контейнерных образов и зависимостей.
    \item Разделять базовую CPU-конфигурацию и профиль аппаратного ускорения через переменные окружения.
    \item Сохранять артефакты прогонов: логи, snapshots Grafana, выгрузки метрик.
    \item Проводить повторные прогоны при неизменной конфигурации для оценки разброса метрик.
\end{enumerate}

\section*{Б.5 Переменные окружения (пример)}
\phantomsection
\addcontentsline{toc}{section}{Б.5 Переменные окружения (пример)}

\begin{longtable}{@{}p{0.34\textwidth}p{0.6\textwidth}@{}}
\nolinkurl{POSTGRES_DB} & Имя базы данных платформы \\
\nolinkurl{POSTGRES_USER} & Пользователь БД \\
\nolinkurl{POSTGRES_PASSWORD} & Пароль пользователя БД \\
\nolinkurl{DATABASE_URL} & Строка подключения API к PostgreSQL \\
\nolinkurl{RABBITMQ_URL} & Строка подключения к RabbitMQ \\
\nolinkurl{JWT_SECRET} & Секрет подписи JWT \\
\nolinkurl{JWT_ACCESS_TTL_MIN} & Время жизни access-токена (минуты) \\
\nolinkurl{JWT_REFRESH_TTL_DAYS} & Время жизни refresh-токена (дни) \\
\nolinkurl{WEBHOOK_TARGET_URL} & URL целевого HTTP-получателя (в тестах - \texttt{webhook-mock}) \\
\nolinkurl{OUTBOX_BATCH_SIZE} & Размер пакета публикации таблицы исходящих сообщений \\
\nolinkurl{OUTBOX_POLL_INTERVAL_MS} & Интервал опроса таблицы исходящих сообщений \\
\nolinkurl{DELIVERY_MAX_RETRIES} & Лимит повторов доставки \\
\nolinkurl{DELIVERY_BACKOFF_BASE_MS} & Базовая задержка между повторами доставки \\
\nolinkurl{PROMETHEUS_ENABLED} & Флаг экспорта метрик \\
\end{longtable}

\section*{Б.6 Фрагмент docker-compose (концептуально)}
\phantomsection
\addcontentsline{toc}{section}{Б.6 Фрагмент docker-compose (концептуально)}

\begin{verbatim}
services:
  api:
    build: ./services/api
    depends_on: [postgres, rabbitmq]
  analytics-worker:
    build: ./services/analytics-worker
    depends_on: [postgres, rabbitmq]
  integration-worker:
    build: ./services/integration-worker
    depends_on: [postgres, rabbitmq, webhook-mock]
  ui:
    build: ./services/ui
    depends_on: [api]
  postgres:
    image: postgres:16
  rabbitmq:
    image: rabbitmq:3-management
  prometheus:
    image: prom/prometheus
  grafana:
    image: grafana/grafana
  webhook-mock:
    image: mendhak/http-https-echo
\end{verbatim}

\section*{Б.7 Профиль дальнейшего аппаратного ускорения}
\phantomsection
\addcontentsline{toc}{section}{Б.7 Профиль дальнейшего аппаратного ускорения}

При переносе на стенд с аппаратным ускорением рекомендуется:
\begin{enumerate}
    \item Включить профиль контейнеров с доступом к ускорителю.
    \item Отдельно зафиксировать версию драйвера NVIDIA и среды исполнения.
    \item Повторно откалибровать \(\tau_{\text{conf}}\), размер входа модели и целевой FPS.
    \item Выполнить серию прогонов 1, 3, 5, 10 камер с неизменными правилами.
    \item Для каждого шага мощности зафиксировать p95 \texttt{object->event}, dropped frames и delivery failures.
\end{enumerate}

Эти шаги необходимы для сопоставимости профилей запуска и для корректной интерпретации масштабируемости платформы.
